\chapter{Logical View}
\ChapterLabel{LogicalView}
\label{Section:LogicalViewResponsibility}
This chapter defines conceptual responsibilities and data boundaries without selecting C++ classes.


\section{Responsibility and Contract Map}
These are conceptual responsibilities, not implemented classes. The register identifies their inputs, outputs, and governing requirements. No class diagram or C++ API is selected.

\small
\begin{tabularx}{\linewidth}{L{3.0cm}Y}
\caption{Measurement Responsibilities\label{Table:MeasurementResponsibilities}}\\
\toprule
Responsibility & Inputs, outputs, and governing links\\\midrule
\endfirsthead
\toprule
Responsibility & Inputs, outputs, and governing links\\\midrule
\endhead
\bottomrule
\endlastfoot
Analog conversion & \RequirementRef{REQ-F-101}; sensor signal to ADC0804 8-bit conversion code. Analog span/reference/calibration: \hyperref[Question:Q-102]{Q-102}.\\
Counter observation & \RequirementRef{REQ-F-104}; clock/counter to counter value. Width, event and rollover: \hyperref[Question:Q-111]{Q-111}.\\
Capture and association & \RequirementRef{REQ-F-105}; conversion and counter inputs to an associated acquisition record. No record layout is selected: \hyperref[Question:Q-112]{Q-112}.\\
FT245R transfer and host reception & \RequirementRef{REQ-F-102}, \RequirementRef{REQ-F-103}; associated acquisition data across parallel/USB boundaries. Framing and partial records: \hyperref[Question:Q-103]{Q-103}.\\
Qualification & \RequirementRef{REQ-F-106}; received measurement evidence to a usability disposition under \hyperref[Question:Q-113]{Q-113}.\\
Publication & \RequirementRef{REQ-F-107}, \RequirementRef{REQ-I-101}; code, status, timing to downstream consumer under \hyperref[Question:Q-114]{Q-114}.\\
Source condition & \RequirementRef{REQ-F-108}, \RequirementRef{REQ-F-109}; defined observations to Field-Engineer connection/health information under \hyperref[Question:Q-115]{Q-115}.\\
iOS presentation & \RequirementRef{REQ-C-104}, \RequirementRef{REQ-F-110}; processing-core reading and status outputs to the iOS SwiftUI UI. The UI/core contract remains \hyperref[Question:Q-106]{Q-106}; presentation semantics remain \hyperref[Question:Q-105]{Q-105}.\\
\end{tabularx}
\normalsize

\section{Data Meaning and Invariants Awaiting Definition}
\label{Section:DataMeaning}
Conversion code, counter value, host qualification status, and published code are separate terms. The published code is not assumed to be the raw conversion code. Counter value is not assumed to be elapsed time, absolute time, or the host receive timestamp. A record's association, timing origin, field units, lifecycle, and representation are defined only after \hyperref[Question:Q-111]{Q-111}, \hyperref[Question:Q-112]{Q-112}, \hyperref[Question:Q-114]{Q-114} is resolved. No packed structure, byte order, checksum, or wall-clock mapping is introduced here.
